\documentclass[12pt,a4paper]{article}
\usepackage[T1]{fontenc}
\usepackage[utf8]{inputenc}
\usepackage{lmodern}
\usepackage[english]{babel}
\usepackage[margin=2.5cm]{geometry}
\usepackage{xcolor}
\usepackage{titlesec}
\usepackage{enumitem}
\usepackage{caption}
\usepackage{microtype}
\usepackage{fancyvrb}
\usepackage{float}
\newfloat{listing}{H}{lol}
\floatname{listing}{Listing}
% Shared by report.tex and presentation.tex.
\usepackage{graphicx}
\usepackage{booktabs}
\usepackage{tikz}
\usetikzlibrary{arrows.meta,positioning,fit,calc,shapes.misc}

\definecolor{todored}{HTML}{C0392B}
\definecolor{good}{HTML}{1E8449}
\definecolor{bad}{HTML}{B03A2E}
\definecolor{mesh}{HTML}{2E86C1}

% Anything the team still has to fill in. Search the .tex files for \todo.
\newcommand{\todo}[1]{{\color{todored}\textbf{[TODO:} #1\textbf{]}}}

% \placeholder{width}{height}{text}: red box where a figure will go.
\newcommand{\placeholder}[3]{%
  \fcolorbox{todored}{todored!5}{\parbox[c][#2][c]{\dimexpr#1-2\fboxsep-2\fboxrule\relax}%
    {\centering\color{todored}\small #3}}}

% \evidence{file}{width}{height}{what it should show}
% Shows evidence/<file> once it exists, a placeholder until then.
\newcommand{\evidence}[4]{%
  \IfFileExists{evidence/#1}%
    {\includegraphics[width=#2,height=#3,keepaspectratio]{evidence/#1}}%
    {\placeholder{#2}{#3}{missing \texttt{evidence/#1}\\[2pt]#4}}}

\tikzset{
  >={Stealth},
  svc/.style={draw, rounded corners, minimum width=2.3cm, minimum height=.75cm,
              font=\small\ttfamily, fill=white},
  hacked/.style={svc, draw=bad, fill=bad!12},
  proxy/.style={svc, fill=mesh!15, draw=mesh},
}

\usepackage[hidelinks]{hyperref}

\titlelabel{\thetitle.\quad}
\setlength{\parskip}{0.5em}
\setlength{\parindent}{0pt}

\begin{document}

% ---------------------------------------------------------------- title page
\begin{titlepage}
  \centering
  {\large Lodz University of Technology\\
  Faculty of Technical Physics, Information Technology and Applied Mathematics\\
  Institute of Information Technology\par}
  \vspace{1.5cm}
  {\large Selected Issues of Network Applications\\2026/2027\par}
  \vspace{3cm}
  {\LARGE\bfseries Modern Service Mesh: a better way to secure communication
  between microservices\par}
  \vspace{1cm}
  {\Large Stage 3 report: final results\par}
  \vspace{3cm}
  {\large
  Jędrzej Bartoszewski, 251482\\
  Michał Domagała, 251505\\
  Kacper Maziarz, 251586\\
  Mikołaj Pawłoś, 258681\par}
  \vfill
\end{titlepage}

The problem definition, the choice of tools and the functional requirements
are unchanged from the stage 1 report, and steps 1 and 2 are described in the
stage 2 report. This report covers steps 3--5: the authorization policies, the
monitoring dashboard and the attack simulation, and what we learned while
building them.

Everything runs on a local k3d cluster with one server and two agent nodes.
The setup scripts pin k3d v5.9.0, kubectl v1.37.2, Linkerd edge-26.11.1 and
Gateway API v1.2.1, and the manifests pin Online Boutique v0.10.7 and Grafana
11.3.0. The scripts, manifests and policies named below are in the Project
Resources archive submitted with this report; paths are relative to its
root.

% ---------------------------------------------------------------- 1
\section{Work progress}

\begin{table}[H]
  \centering
  \small
  \begin{tabular}{cllll}
    \toprule
    \# & Step & Planned & Owner & Status \\
    \midrule
    1 & k3d cluster and Online Boutique      & Weeks 3--4                  & Mikołaj  & Done \\
    2 & Linkerd with mTLS (Proof of Concept) & Weeks 5--6, until 6.11.2026 & Jędrzej  & Done \\
    3 & Zero-Trust authorization policies    & Week 7                      & Michał   & Done \\
    4 & Prometheus and Grafana               & Week 8                      & Kacper   & Done \\
    5 & Attack simulation                    & Week 9, until 27.11.2026    & Michał   & Done \\
    \bottomrule
  \end{tabular}
  \caption{Schedule and work division from the stage 1 report, with the final
    status.}
\end{table}

\subsection*{Step 3: Zero-Trust authorization policies}

After step 2 every connection in the mesh was encrypted and both sides knew
each other's identity, but every service still accepted requests from anyone.
Step 3 adds authorization: each service now accepts only the callers and the
operations that Online Boutique needs, and refuses everything else.

\textbf{Allow-matrix.} We first recorded which identity calls which gRPC
method during normal use: about 90 seconds of traffic from
\texttt{loadgenerator}, the shop's built-in component that simulates
customers, plus three orders placed by hand. We captured it with
\texttt{linkerd viz tap}, which streams the requests passing through the
proxies. Of the 16 connections between services seen in step 2, 14 use gRPC,
and they broke down into 20 distinct (caller, method) pairs; the other two
are \texttt{loadgenerator} to \texttt{frontend} (HTTP) and
\texttt{cartservice} to \texttt{redis-cart} (TCP). Table~\ref{tab:matrix}
is the result. The
callers are Kubernetes ServiceAccounts: Online Boutique gives each service its
own, and Linkerd puts it into the service's mTLS certificate as its identity.

\begin{table}[H]
  \centering
  \small
  \begin{tabular}{lll}
    \toprule
    Service & Methods & Allowed callers \\
    \midrule
    \texttt{paymentservice}        & \texttt{Charge}                       & checkout \\
    \texttt{emailservice}          & \texttt{SendOrderConfirmation}        & checkout \\
    \texttt{checkoutservice}       & \texttt{PlaceOrder}                   & frontend \\
    \texttt{adservice}             & \texttt{GetAds}                       & frontend \\
    \texttt{recommendationservice} & \texttt{ListRecommendations}          & frontend \\
    \texttt{cartservice}           & \texttt{GetCart}, \texttt{EmptyCart}  & frontend, checkout \\
                                   & \texttt{AddItem}                      & frontend \\
    \texttt{currencyservice}       & \texttt{Convert}                      & frontend, checkout \\
                                   & \texttt{GetSupportedCurrencies}       & frontend \\
    \texttt{shippingservice}       & \texttt{GetQuote}                     & frontend, checkout \\
                                   & \texttt{ShipOrder}                    & checkout \\
    \texttt{productcatalogservice} & \texttt{GetProduct}                   & frontend, checkout \\
                                   & \texttt{ListProducts}                 & frontend, recommendation \\
                                   & \texttt{SearchProducts}               & nobody \\
    all nine above                 & gRPC health check                     & kubelet (node addresses) \\
    \texttt{frontend}              & everything (HTTP)                     & any network \\
    \texttt{redis-cart}            & whole port (TCP)                      & cart; kubelet \\
    \bottomrule
  \end{tabular}
  \caption{Who may call what after step 3. Callers are written without the
    \texttt{service} suffix: ``checkout'' is \texttt{checkoutservice}.
    \texttt{SearchProducts} exists in the API, but nothing in the shop calls
    it. The kubelet, the node agent that runs the health checks, has no mesh
    identity, so it is recognised by the address its checks come from. Each
    node has its own pod network (\texttt{10.42.0.0/24},
    \texttt{10.42.1.0/24}, \texttt{10.42.2.0/24}), and the checks come from
    its gateway address (\texttt{10.42.0.1} and so on).
    \texttt{loadgenerator} receives no traffic and gets no policy.}
  \label{tab:matrix}
\end{table}

\textbf{Policies.} Linkerd describes authorization with a few Kubernetes
resources that refer to each other. A \texttt{Server} selects the pods of one
service and one of their ports. A \texttt{GRPCRoute} attached to the
\texttt{Server} matches a set of gRPC methods. An
\texttt{AuthorizationPolicy} attached to a route says who may use it, by
pointing to a \texttt{MeshTLSAuthentication} (a list of allowed identities)
or a \texttt{NetworkAuthentication} (a list of allowed source addresses).
A request reaches the application only if it matches a route and passes that
route's policy.

The matrix is written this way in \texttt{policies/}, one file per service:
11 \texttt{Server}s, 22 \texttt{GRPCRoute}s, 25 \texttt{AuthorizationPolicy}
objects, and 7 shared authentications. Routes group the methods that have the
same callers, so the policies work per method (layer 7): \texttt{frontend}
may call \texttt{GetQuote} on \texttt{shippingservice} but not
\texttt{ShipOrder}. The two services that don't speak gRPC are authorized per
port: \texttt{frontend} is the public entry point and accepts anyone, and
\texttt{redis-cart} accepts only \texttt{cartservice}.
Listing~\ref{lst:payment} shows the rule for \texttt{Charge}.


\textbf{Rollout.} We applied the policies first in audit mode, in which the
proxy lets a request through but records that the policies would have denied
it, both in its metrics and in its log. We then placed an order through the
shop and ran the attacker pods described below. On the nine gRPC services
and on \texttt{redis-cart}, every audited request came from an attacker pod
and none from a legitimate caller, so we switched to enforcing. The
\texttt{frontend} reports
differently: its policy is attached to the whole \texttt{Server} rather than
to a route, and in audit mode the proxy lists all such traffic as audited,
including the loadgenerator's, even though the policy allows it. \texttt{./zero-trust.sh on} and \texttt{off} (and the PowerShell
version) switch the policies without restarting any pod, which is what the
before/after demo in step 5 needs. \texttt{setup.sh -{}-protect} builds a new
cluster with the policies and the namespace default-deny (below) already on.

\begin{listing}[H]
\begin{Verbatim}[fontsize=\footnotesize,frame=single,framesep=2mm]
kind: Server                                  kind: AuthorizationPolicy
metadata: {name: paymentservice}              metadata: {name: paymentservice-charge}
spec:                                         spec:
  podSelector:                                  targetRef:
    matchLabels: {app: paymentservice}            kind: GRPCRoute
  port: 50051                                     name: paymentservice-charge
  proxyProtocol: gRPC                           requiredAuthenticationRefs:
---                                             - kind: MeshTLSAuthentication
kind: GRPCRoute                                   name: callers-checkout
metadata: {name: paymentservice-charge}       ---
spec:                                         kind: MeshTLSAuthentication
  parentRefs:                                 metadata: {name: callers-checkout}
  - {kind: Server, name: paymentservice}      spec:
  rules:                                        identityRefs:
  - matches:                                    - kind: ServiceAccount
    - method:                                     name: checkoutservice
        service: hipstershop.PaymentService
        method: Charge
\end{Verbatim}
\caption{The four resources that let only \texttt{checkoutservice} call
  \texttt{PaymentService/Charge}, shortened (no \texttt{apiVersion},
  namespace or labels). The same file has a second route and policy for the
  health check. Full version: \texttt{policies/paymentservice.yaml} and
  \texttt{policies/00-authentications.yaml}.}
\label{lst:payment}
\end{listing}

\textbf{Verification.} Three test pods act as attackers: one outside the mesh
(no identity), one running as \texttt{adservice} (a compromised minor
service) and one running as \texttt{frontend} (a compromised service that has
some legitimate access). \texttt{attack/verify.sh} uses them to check the
whole matrix in 47 cases (Table~\ref{tab:verify}). Each gRPC call carries an
empty message, except \texttt{ShipOrder}, which gets an empty address because
a completely empty request crashes \texttt{shippingservice}
(section~\ref{sec:insights}). That is enough to tell the outcomes apart. The
report uses four gRPC status codes: 0 (\texttt{OK}) and 2
(\texttt{UNKNOWN}, here an error raised by the application) mean the call
reached the application; 7 (\texttt{PERMISSION\_DENIED}) means the proxy
refused it; 5 (\texttt{NOT\_FOUND}) is what the proxy returns when a request
matches no route.

\begin{table}[H]
  \centering
  \footnotesize
  \begin{tabular}{p{8.6cm}rl}
    \toprule
    Check & Cases & Result with the policies on \\
    \midrule
    every method in the matrix, from the outside and \texttt{adservice} pods
      & 28 & denied (status 7) \\
    \texttt{Charge}, \texttt{ShipOrder}, \texttt{SendOrderConfirmation} from
      the \texttt{frontend} pod & 3 & denied \\
    \texttt{GetQuote} from the \texttt{frontend} pod & 1 & reaches the application \\
    \texttt{SearchProducts}, from all three pods & 3 & no route (status 5) \\
    health check of \texttt{paymentservice}, from all three pods & 3 & denied \\
    Redis \texttt{PING}, from all three pods & 3 & no answer \\
    shop: home page, add to cart, checkout & 3 & order confirmed \\
    pods \texttt{2/2} and no restarts; no denial counted for an allowed caller
      & 3 & as expected \\
    \midrule
    total & 47 & all passed \\
    \bottomrule
  \end{tabular}
  \caption{What \texttt{attack/verify.sh} checks. Redis speaks plain TCP, so
    a refused connection is simply closed and the attacker gets no reply.
    \texttt{2/2} means each pod runs both its application container and the
    Linkerd proxy. With the policies off, \texttt{verify.sh -{}-expect-open}
    runs the 43 cases that apply there, and every attack case reaches the
    application.}
  \label{tab:verify}
\end{table}

\textbf{Namespace default-deny.} The policies cover every application port,
but a port without a \texttt{Server} still falls back to Linkerd's default,
\texttt{all-unauthenticated}. That leaves the proxy admin port 4191 open on
every pod, so any pod can read another service's traffic metrics. Switching
the namespace default to \texttt{deny} closes it, together with one extra
policy that lets only the linkerd-viz Prometheus read \texttt{/metrics} and
only the kubelet reach the proxy's own health checks
(\texttt{policies/namespace-default-deny/}, \texttt{./zero-trust.sh deny
on}). We tested it: \texttt{verify.sh} still passed, Prometheus kept
scraping, and an attacker pod got HTTP 403 on \texttt{/metrics}. It is part
of the
protected configuration: \texttt{setup.sh -{}-protect} sets the default before
the shop starts, so the pods start with it and nothing has to restart. On a
running cluster the switch is slower, because the proxy reads the default only
when it starts, so turning it on or off restarts every pod. With \texttt{deny}
in place, removing the policies also no longer opens the shop but breaks it.
The live before/after demo therefore uses a cluster set up without
\texttt{-{}-protect} and switches only the per-service policies, which takes
seconds.

\subsection*{Step 4: Prometheus and Grafana}

Grafana runs in its own namespace, inside the mesh, and reads the Prometheus
that comes with linkerd-viz, so nothing else has to collect metrics. Of the
two metric stores considered in stage 1, Prometheus and VictoriaMetrics, this
made Prometheus the obvious choice: it was already running and scraping every
proxy. The setup
scripts install linkerd-viz before the shop, so every application pod starts
ready for \texttt{linkerd viz tap}, and then apply \texttt{observability/}
with \texttt{kubectl apply -k}. The dashboard ``Zero Trust: allowed vs
denied'' opens at \url{http://localhost:3000} without logging in
(Figure~\ref{fig:grafana}).

Its panels are built from three counters that every Linkerd proxy exports:
requests allowed by a policy, requests denied, and TCP connections denied.
Prometheus scrapes them every 10 seconds and the dashboard refreshes every 5,
so an attack shows up within about 15 seconds. The top row has the totals:
denied calls since the pods started, allowed requests per second, the share
of allowed calls that came with a verified mesh identity, and denied TCP
connections and denied health checks over the last five minutes. Denied
health checks must stay at zero, otherwise the policies are breaking the
pods. Below are the denied and allowed rates per service,
and the table ``Who was denied, and where'' with the target service, the
route and the caller of every refusal. A caller outside the mesh has no
identity and is shown with the reason \texttt{no\_tls\_from\_remote}.
The share of allowed calls with a mesh identity stays just under 100\,\%
because the \texttt{frontend} also receives traffic from outside the mesh:
the shop's visitors through the k3d load balancer and the kubelet's health
checks.

We tested the dashboard twice. On a cluster built from scratch on Windows,
three runs of the attack script gave 12 refused calls and the dashboard
counted 12. A second run on an existing cluster gave the same result, with no
denied health checks and no denied TCP connections; this time every dashboard
query was also run through Grafana's API, with the same numbers.

\begin{figure}[H]
  \centering
  \evidence{10-grafana.png}{\linewidth}{9cm}{Grafana dashboard right after
    \texttt{attack/attack.sh}}

  \vspace{0.5em}
  \evidence{10b-grafana-denied.png}{\linewidth}{3cm}{the ``Who was denied''
    table, enlarged}
  \caption{The dashboard after the attack runs with the policies on: 16
    denied calls, all on \texttt{paymentservice} and
    \texttt{cartservice}, half from the \texttt{adservice} identity and half
    from the pod outside the mesh. The gap in the allowed-requests chart
    around 19:48 is the run with the policies off: the panels count only
    traffic that passes through our \texttt{Server}s, and without the
    policies there are none, so in the ``before'' state the dashboard shows
    nothing at all. Bottom: the table of refusals, enlarged.}
  \label{fig:grafana}
\end{figure}

\subsection*{Step 5: attack simulation}

The attack repeats the lateral movement from the stage 2 report, where we
showed that mTLS alone doesn't stop a compromised pod from calling the
payment service. Two of the attacker pods, the one outside the mesh and the
one with the \texttt{adservice} identity, each call
\texttt{PaymentService/Charge} and \texttt{CartService/GetCart} directly,
four calls in total (\texttt{attack/attack.sh}). The script is the short
version of \texttt{verify.sh}, meant for the live demo.

Table~\ref{tab:attack} shows the two runs. Without the policies
(\texttt{./zero-trust.sh off}) all four calls reach the application.
\texttt{Charge} comes back with the payment service's own error about the
missing card number, and \texttt{GetCart} succeeds. With the policies on
(\texttt{./zero-trust.sh on}) all four end with \texttt{PermissionDenied}
before reaching the application, and the Grafana dashboard lists each refusal
with the caller's identity (Figure~\ref{fig:grafana}). That the shop works
in both states is checked by \texttt{verify.sh}, run once with the policies
off and once with them on (Table~\ref{tab:verify}).

\begin{table}[H]
  \centering
  \small
  \begin{tabular}{llll}
    \toprule
    Attacker pod & Call & Policies off & Policies on \\
    \midrule
    outside the mesh   & \texttt{Charge}  & status 2, ``reading \texttt{credit\_card\_number}'' & status 7 \\
                       & \texttt{GetCart} & status 0 (success)                               & status 7 \\
    \texttt{adservice} & \texttt{Charge}  & status 2, ``reading \texttt{credit\_card\_number}'' & status 7 \\
                       & \texttt{GetCart} & status 0 (success)                               & status 7 \\
    \bottomrule
  \end{tabular}
  \captionsetup{justification=raggedright}
  \caption{gRPC status of each call made by \texttt{attack/attack.sh}. Status 2
    is the payment service failing on the empty card data, which means the
    call reached it; status 7 is \texttt{PermissionDenied} from the proxy.}
  \label{tab:attack}
\end{table}

% ---------------------------------------------------------------- 2
\section{Insights and challenges}
\label{sec:insights}

\subsection*{Routes are mandatory, so per-method rules cost little}

We expected method-level rules to be much more work than port-level ones. They
weren't, because of how Linkerd behaves once a port has a \texttt{Server}.
A \texttt{Server} with no authorization denies everything on its port,
including the kubelet's health checks, and the pod restarts. The health
checks have to be allowed by address, because the kubelet has no identity,
while application calls are allowed by identity. Those are two different
policies, so even a port-level setup needs a separate route for the health
check. And once a port has any route, a request that matches no route is
rejected with gRPC status 5 (``no route found''), so the application calls
need a route too. Splitting that second route into one route per method group
adds little. It also means that a method nothing uses, like
\texttt{SearchProducts}, gets no route and can't be called at all.

Two smaller rules shaped the files. An \texttt{AuthorizationPolicy} can list
several required authentications, but all of them must match, so ``frontend
or checkout'' has to be a single \texttt{MeshTLSAuthentication} that lists
both identities. And the health checks come from the node's gateway address,
one per node, so the policies list the three addresses of our three nodes.
That ties the policies to the cluster layout. It hit us once: after an
interrupted setup run, one node got a different pod network, and
\texttt{checkoutservice}, \texttt{emailservice} and
\texttt{recommendationservice} kept failing their health checks and
restarting, and checkout stopped working.

\subsection*{What least privilege does for a compromised service}

The \texttt{attacker-as-frontend} pod shows what the policies can't do: a pod
with the frontend's identity still has the frontend's rights. It can read
shipping quotes and carts, because the real frontend must. It can no longer
call \texttt{Charge}, \texttt{SendOrderConfirmation} or \texttt{ShipOrder}
directly, which it could do before step 3. It can still trigger all three by
placing an order through \texttt{checkoutservice}, because the real frontend
must be able to do that too; the policies decide who may call what, not what
a permitted call leads to. Without a service mesh, the only way to express this would be in the
application code of every service. \texttt{NetworkPolicy}, the traffic filter
built into Kubernetes, works on IP addresses and labels and doesn't see gRPC
methods at all.

The test pods also turned up a weak spot in the application itself: an empty
\texttt{ShipOrder} request crashes \texttt{shippingservice}, because it reads
the address without checking that it is there. Without policies any pod in the
cluster can restart that service at will. With them, only
\texttt{checkoutservice} can send \texttt{ShipOrder}.

\subsection*{Switching policies on without breaking the shop}

A \texttt{Server} that exists without its routes and policies denies every
request on its port. That can happen for a moment in both directions. When
the policies are applied, we put each \texttt{Server} last in its file, so it
is created after its routes and policies. When they are removed, we delete
them with one command that selects them by label and lists the
\texttt{Server}s as the first resource type, so they go before the rest.
Removal now loses no request. Applying can still refuse a few requests,
because the proxies receive the new rules a moment apart and a
\texttt{Server} can arrive before its routes: in three remove/apply cycles
this hit 0 to 5 \texttt{frontend} requests per apply and an occasional
health check, which does not restart a pod (\texttt{policies/README.md}).
For a demo we accept that.

\subsection*{Testing the tests}

Our first version of \texttt{verify.sh} could pass on a broken setup. It
treated any status other than 7 as ``reached the application'', because the
application's own errors were the sign of a call getting through. A missing
route returns status 5, so a route we had forgotten looked like success. An
attacker pod that started without its proxy looked like the outside
attacker. And the script never checked whether legitimate callers were being
refused. We broke the cluster on purpose for each case, watched the script
fail, then fixed it. The current version also checks the attacker pods'
identities and compares the proxies' denial counters for allowed callers
before and after the run.

\subsection*{Running it all on a laptop}

With Linkerd, linkerd-viz and Grafana, the cluster needs 6--7~GB of memory in
Docker. With 4~GB, the Docker VM started swapping and the cluster stopped
responding, twice. Overload has a side effect that is easy to miss. If the
Linkerd proxy injector is down while a pod starts, the pod starts without a
proxy (\texttt{1/1} instead of \texttt{2/2}). It works, but it has no
identity, so the policies treat it as an outsider and the shop partly breaks.
The setup scripts now detect such pods and recreate them.

A cluster created before step 4 also has no port 3000 on its load balancer,
because k3d sets the ports only when it creates the cluster. Grafana then runs
but can't be reached from the host until the port is added with
\texttt{k3d cluster edit}.

\subsection*{Monitoring needs its own authorization}

linkerd-viz protects its own Prometheus: only its own \texttt{metrics-api}
may query it. Without an extra policy every Grafana query fails with HTTP 403,
and the Prometheus proxy counts the denial for the Grafana identity. We had to put Grafana in the
mesh, so that it has an identity, and add a policy that lets that identity
query Prometheus. The policies protect the monitoring stack just as they
protect the shop.

The namespace default-deny adds a second effect. On a cluster set up with
\texttt{-{}-protect}, the proxies refused between 18 and 32 requests for
\texttt{/metrics} per pod in the first minutes after the shop started, all
sent without mTLS, while Prometheus was apparently not yet treating the new
pods as meshed. After that every scrape used mTLS and was allowed, so the
metrics are complete once the shop has settled, but the dashboard counts
these refusals as denied calls.

The per-second chart also misses the first denial of a new caller and method
pair, because that denial creates a new counter and a rate needs two samples.
The ``since pod start'' total and the table do count it, so for the demo we
run the attack more than once.

\subsection*{Resource cost of the mesh and the policies}

We measured the overhead with \texttt{fortio} in three variants: the shop
without the mesh, with mTLS only, and with mTLS and the policies
(Table~\ref{tab:perf}). Each number is the median of three 60-second runs at
a fixed 20 requests per second, with the shop redeployed when the mesh was
added or removed.

\begin{table}[H]
  \centering
  \small
  \begin{tabular}{lrrr}
    \toprule
    Median latency (p50) & No mesh & mTLS & mTLS + policies \\
    \midrule
    \texttt{ListProducts}, direct gRPC call & 0.8 ms  & 1.7 ms  & 1.7 ms \\
    product page (\texttt{frontend})       & 10.9 ms & 22.7 ms & 22.6 ms \\
    home page (\texttt{frontend})          & 16.4 ms & 28.3 ms & 27.4 ms \\
    \bottomrule
  \end{tabular}
  \caption{Latency overhead of the mesh and the policies. Details, CPU,
    throughput and the raw runs: \texttt{stage4/performance/}.}
  \label{tab:perf}
\end{table}

The mesh (a proxy on each side of a call, with mTLS between them) roughly
doubles the latency of a single call, adding about 1~ms per hop, which adds up
to about 12~ms on a page that makes many calls. We did not measure the cost of
the encryption separately from the rest of the proxies' work. The policies
add no measurable latency on top of that: the medians differ by less than
1~ms. The maximum throughput of the pages stayed the same in all three
variants (45--47 and 56--58 requests per second), because it is limited by
the shop's own CPU limits. In the page tests, the proxies of all shop pods
together used about 280 millicores of CPU at 20 requests per second; in the
direct gRPC test, about 27 millicores.

Proxy memory readings differed between measurements. In the performance runs
the proxies used about 3.6~MiB each, as in stage 2, and 3--6~MiB on the same
cluster more than two hours after the shop was redeployed. Separate
\texttt{kubectl top} snapshots showed 12--16~MiB per proxy, the same with the
policies on and off. We did not investigate why. The performance numbers come
from one laptop and only support comparisons between the variants.

\subsection*{Limits of the delivered solution}

\begin{itemize}[itemsep=2pt]
  \item On a cluster without the namespace default-deny, such as the one for
    the live before/after demo, the proxy admin port 4191 stays open on
    every pod, and any pod can read another service's traffic metrics.
  \item The health checks are allowed by node address, so the policies assume
    exactly three nodes with the pod networks k3d gave them. A different
    cluster needs those addresses changed.
  \item A compromised service keeps its own rights. The policies limit what
    it can reach, not what it does with the access it legitimately has.
  \item The Prometheus in linkerd-viz keeps six hours of data and loses it
    when its pod restarts, and the counters start from zero when a proxy
    restarts. That is enough for a demo, not for an audit trail.
  \item Denied TCP connections to \texttt{redis-cart} carry no method, so the
    dashboard can't tell an attacker from a misconfigured health check there.
  \item Everything was tested on one Linkerd edge release, edge-26.11.1. Edge
    releases can change behaviour between versions.
\end{itemize}

% ---------------------------------------------------------------- 3
\section{Summary}

After step 2, the mesh encrypted every connection and gave each service a
verifiable identity, but a compromised pod could still call the payment
service directly. The authorization policies from step 3 close that gap.
Every application port now accepts only the callers and methods the shop
uses, and with the namespace default-deny, which the protected setup switches
on, so does the proxy's admin port. The attack from stage 2 ends with \texttt{PermissionDenied}, and the
shop keeps working. None of this needed a change to the application. The Grafana
dashboard from step 4 shows each refusal and who caused it.

The project could be taken further in a few directions that follow from the
limits above: a separate Prometheus with longer retention and alerts on
denied requests, health-check rules that don't depend on node addresses, and
running \texttt{verify.sh} automatically after every change to the policies,
so that a missed caller shows up before it breaks the shop.


\begin{table}[H]
  \centering
  \footnotesize
  \begin{tabular}{p{5.2cm}p{8.8cm}}
    \toprule
    Requirement (stage 1) & How it is met \\
    \midrule
    Automated mTLS encryption & Linkerd proxies in every pod; all 16
      connections between services secured (stage 2). \\
    Service identity verification & every proxy has a certificate for its
      ServiceAccount; the policies authorize by that identity. \\
    Access control & 11 \texttt{Server}s: per-method routes and policies on
      the nine gRPC services, per-port rules for \texttt{frontend} and
      \texttt{redis-cart}, namespace default-deny for everything else (step 3). \\
    Observability & Grafana dashboard of allowed and denied traffic on the
      linkerd-viz Prometheus (step 4). \\
    Security validation & packet capture without and with mTLS (stage 2); the
      same attack with mTLS only and with the policies, and 47 cases in
      \texttt{verify.sh} (step 5). \\
    \bottomrule
  \end{tabular}
  \caption{The functional requirements from the stage 1 report.}
\end{table}

\end{document}
